\documentclass[11pt]{article}
\usepackage[T1]{fontenc}
\usepackage{amsmath,amssymb}
\usepackage{cite}
\usepackage{booktabs}
\usepackage{graphicx}
\usepackage{xurl}
\usepackage{hyperref}
\hypersetup{
  colorlinks=true,
  linkcolor=blue,
  urlcolor=blue,
  citecolor=blue,
  pdfauthor={Jabulani Monchwe, Martin Chanza, Olorato Motitswane, Mfundo Monchwe}
}
\newcommand{\software}{\textsf{ReliableML}}
\newcommand{\code}[1]{\texttt{\small #1}}
\newcommand{\jabulaniauthor}{
\IEEEauthorblockN{Jabulani Monchwe}
\IEEEauthorblockA{\textit{Applied Business Analytics and Decision Making}\\
\textit{North-West University}\\
\textit{IDEA-Research Club Chairperson}\\
Johannesburg, South Africa\\
\href{mailto:monchwej.25@gmail.com}{monchwej.25@gmail.com}}
}
\newcommand{\martinchanzauthor}{
\IEEEauthorblockN{Martin Chanza}
\IEEEauthorblockA{\textit{Applied Business Analytics and Decision Making}\\
\textit{North-West University}\\
Mafikeng, South Africa\\
\href{mailto:martin.chanza@nwu.ac.za}{martin.chanza@nwu.ac.za}}
}
\newcommand{\oloratoauthor}{
\IEEEauthorblockN{Olorato Motitswane}
\IEEEauthorblockA{\textit{Eduvos Midrand Campus -- IT Faculty Data Science, South Africa}\\
\textit{IDEA-Research Club}\\
\href{mailto:oloratomotitswane@gmail.com}{oloratomotitswane@gmail.com}}
}
\newcommand{\mfundoauthor}{
\IEEEauthorblockN{Mfundo Monchwe}
\IEEEauthorblockA{\textit{Eduvos Midrand Campus -- IT Faculty Data Science, South Africa}\\
\textit{IDEA-Research Club}\\
\href{mailto:diditmfundo@gmail.com}{diditmfundo@gmail.com}}
}
\newcommand{\minenhleauthor}{
\IEEEauthorblockN{Minenhle Mchunu}
\IEEEauthorblockA{\textit{Applied Business Analytics and Decision Making}\\
\textit{Eduvos Midrand Campus -- IT Faculty Data Science, South Africa}\\
\textit{IDEA-Research Club}\\
\href{mailto:Mchunu.minenhle02@gmail.com}{Mchunu.minenhle02@gmail.com}}
}
\newcommand{\projecturl}{\url{https://github.com/Jabualani/ML-Reliabilty}}

\title{\software{} Method Note: Reliable Pipeline and Persistence}
\author{\contactauthors}
\date{Version 0.1.1}
\begin{document}
\maketitle
\begin{abstract}
\code{ReliablePipeline} retains the fitted estimator and a training-feature
reference together. It supports inference-time drift reports, optional
visualizations, labeled evaluation, and persistence.
\end{abstract}

\section{Fit and predict}
At fit time, the pipeline records the feature schema and reference
distributions alongside fitting the estimator. At predict time, input schema
is validated, predictions are produced, and the incoming feature batch is
compared with the stored reference. A returned report can include feature
metrics, aggregate summaries, warnings, and paths to optional plots or JSON
reports. The monitoring options support callers that need to avoid retaining
large reference data or producing plots; disabling monitoring removes those
checks and must be an explicit decision.

\section{Unlabeled versus labeled reports}
Prediction-time distribution monitoring can operate without labels. Accuracy,
confusion counts, MCC, fairness comparisons, and deterioration require true
labels and are computed through the separate evaluation or history APIs.
Distribution similarity is not a substitute for performance or calibration.
The report's similarity score is descriptive, not a probability of
correctness.

\section{Persistence}
\code{save()} serializes the estimator and reference profile together so the
same fitted transformation is used after loading. The profile may include
training feature observations used for plots and bootstrap computations;
artifacts therefore require privacy protection similar to the source data.
Python pickle can execute code when loading an untrusted artifact: load only
files from trusted sources and use compatible dependency versions.

\section{Operational checks}
Use identical feature names, ordering, units, and preprocessing at training
and inference. Validate reliability thresholds against representative
historical and shifted samples. Monitor labeled performance when outcomes
arrive. Avoid interpreting a green or low-drift report as a guarantee of
correct prediction.

\section{Interface}
\begin{verbatim}
model = ReliablePipeline(
    estimator=classifier,
    reliability_required=True,
    plot_dir="reliability_plots",
).fit(X_train, y_train)

predictions, report = model.predict_with_report(X_new)
model.save("reliable-model.pkl")
\end{verbatim}

\section{Further reading}
\begin{thebibliography}{1}
\bibitem{reliability}
\software{} API and operational details:
\url{https://github.com/Jabualani/ML-Reliabilty}.
\end{thebibliography}
\end{document}
